\begin{table}[!t]
\centering
\caption{Split acak vs split bersih (YOLO n, seed 0).}
\label{tab:split_acak}
\footnotesize
\begin{tabular}{lcccccc}
\hline
\textbf{Model} & \textbf{Split latih} & \textbf{Set uji} & \textbf{mAP50} & \textbf{mAP50-95} & \textbf{IK95 mAP50-95} & \textbf{Akurasi diagnosis} \\
\hline
YOLOv8n & split acak & test acak (sendiri) & 0.806 & 0.489 & [0.464, 0.519] & 0.933 \\
YOLOv8n & split acak & test sesi baru* & 0.613 & 0.271 & [0.241, 0.306] & 0.819 \\
YOLOv8n & split bersih (final s0) & test bersih & 0.474 & 0.191 & [0.146, 0.231] & 0.749 \\
YOLOv8n & split bersih (final s0) & test sesi baru* & 0.474 & 0.205 & [0.180, 0.240] & 0.649 \\
YOLO11n & split acak & test acak (sendiri) & 0.822 & 0.481 & [0.456, 0.510] & 0.932 \\
YOLO11n & split acak & test sesi baru* & 0.639 & 0.283 & [0.255, 0.323] & 0.866 \\
YOLO11n & split bersih (final s0) & test bersih & 0.584 & 0.236 & [0.199, 0.279] & 0.741 \\
YOLO11n & split bersih (final s0) & test sesi baru* & 0.586 & 0.263 & [0.234, 0.299] & 0.756 \\
\hline
\end{tabular}
\\[2pt]\parbox{\columnwidth}{\scriptsize *Test sesi baru = satu-satunya set yang sebanding (sama untuk kedua model, tanpa kembaran di train mana pun). Kebocoran split acak (val+test): batas atas 92.0\% (kelompok berantai), batas bawah 76.9\% (kembaran langsung). Satu seed (s0).}
\end{table}
